\paragraph{The favourite--longshot bias in betting markets.} The tendency of bettors to overpay for
unlikely outcomes is one of the oldest anomalies in economics.
\begin{itemize}
  \item \citet{griffith1949} found it in American horse racing, and \citet{ali1977} and
  \citet{hausch1981} documented it on large samples of parimutuel races.
  \item \citet{thalerziemba1988} and \citet{sauer1998} survey the early evidence.
  \item The bias is not universal. \citet{woodland1994} find a reverse bias in baseball money-line
  betting, and its size differs between parimutuel and fixed-odds settings.
  \item \citet{whelan2024} shows that fixed-odds markets generate the bias even with competitive
  bookmakers once bettors are risk-averse and heterogeneous.
  \item \citet{moskowitz2021} uses sports betting contracts as a laboratory for asset pricing and
  finds momentum consistent with delayed overreaction.
\end{itemize}

\paragraph{Explanations.} \citet{ottaviani2008} organise the explanations into four groups.
\begin{itemize}
  \item \emph{Risk preference.} \citet{ali1977} infers risk-loving bettors; \citet{golec1998}
  argue that bettors love skewness rather than variance.
  \item \emph{Probability weighting.} \citet{jullien2000} find that cumulative prospect theory
  \citep{kahneman1979,tversky1992} fits racetrack odds better than expected utility.
  \citet{snowbergwolfers2010} use exotic bets to separate the two and conclude that misperception
  of probabilities explains the bias better than risk-love.
  \item \emph{Information.} \citet{shin1991,shin1993} model bookmakers who protect themselves
  against insiders by shading longshot odds, yielding an estimable insider share.
  \item \emph{Belief heterogeneity.} \citet{ali1977} and \citet{gandhi2015} show that dispersed
  beliefs among budget-constrained bettors produce the bias even without preference distortions.
  \citet{gandhi2015} find that a minority of noise traders with dispersed beliefs fits US
  racetrack data better than preference-based models.
\end{itemize}
Our setting can tell these explanations apart in ways racetracks cannot, because we observe who
trades, in which role, at what size and when.

\paragraph{Prediction markets.} Prices in prediction markets are routinely read as probabilities
\citep{wolferszitzewitz2004,wolferszitzewitz2006}, though \citet{manski2006} shows that, with
heterogeneous beliefs, the price need not equal the mean belief.
\begin{itemize}
  \item \citet{pageclemen2013} find that Intrade prices were well calibrated for events
  resolving soon but biased for events far in the future. They show in a model that time to
  expiration itself pushes prices toward the bias.
  \item \citet{rothschildsethi2016} classify trader strategies with account-level Intrade data.
  They document a single trader who accounted for over a third of bets on one candidate and kept
  that candidate's price above a rival exchange's.
  \item On Kalshi, \citet{burgi2025} document a favourite--longshot bias that is stronger in
  prices accepted by takers than by makers. They explain it with a modest overestimation of small
  probabilities.
\end{itemize}
We extend this line of work in three directions:
\begin{itemize}
  \item across the full range of event categories and contract structures;
  \item with all public trades classified by role, size and timing;
  \item with identified on-chain wallets and matched markets on a second venue.
\end{itemize}

\paragraph{Lottery demand and retail trading in financial markets.} The bias has close
analogues in securities markets.
\begin{itemize}
  \item \emph{Lottery-like stocks.} Individual investors overweight stocks with lottery-like
  payoffs \citep{kumar2009}. Such stocks earn low subsequent returns \citep{bali2011}, as
  probability weighting predicts \citep{barberis2008}.
  \item \emph{Options.} Lottery-like equity options earn large negative returns
  \citep{boyer2014}. Retail options traders, now a majority of volume, concentrate in cheap,
  short-dated contracts and lose money on average \citep{bryzgalova2023}.
  \item \emph{Attention.} Attention drives individual investors' buying \citep{barberodean2008}.
  Attention-induced herding by app-based retail traders predicts negative returns
  \citep{barber2022}.
\end{itemize}
A prediction-market longshot is the purest lottery ticket: a cheap claim that pays a large
multiple with small probability, and that is often bought through the same retail apps. Our
results speak to whether lottery demand, and the retail flow behind it, prices such claims in
markets with a known payoff date and no fundamental cash-flow risk.
